\documentclass[10pt,a4paper]{article}
\usepackage[utf8]{inputenc}
\usepackage[T1]{fontenc}
\usepackage[spanish,es-noshorthands]{babel}
\usepackage[margin=1.8cm]{geometry}
\usepackage{booktabs,tabularx,enumitem}
\setlength{\parindent}{0pt}
\setlength{\parskip}{3pt}
\pagestyle{empty}

\begin{document}

{\Large\textbf{Bitácora --- Generador de tablas LR(1)}}\\[2pt]
CS342 Compiladores · UCSP · 2026-II\hfill
Eduardo Cruz Laura · Alexander Carpio Mamani · Fernando Zegarra Maratuech

\medskip
\hrule
\medskip

\textbf{Registro de trabajo}

\smallskip
{\small
\begin{tabularx}{\linewidth}{@{}l l X@{}}
\toprule
\textbf{Fecha} & \textbf{Integrante} & \textbf{Actividad}\\
\midrule
02/09 & Eduardo Cruz & Ajuste de \texttt{emit()} en el scanner para mostrar cada token en consola.\\
29/09 & Eduardo Cruz & Estructura del proyecto \texttt{parser/}: lectura de la gramática desde archivo (alternativas con \texttt{|}, producción vacía \texttt{epsilon}), \texttt{first()} por punto fijo y \texttt{firstCadena($\beta$)}, con sus primeras pruebas.\\
29/09 & Eduardo Cruz & CMake multiplataforma (macOS y Windows) y script de limpieza de builds.\\
29/09 & Fernando Zegarra & Integración de la rama \texttt{prueba-edcr} en \texttt{main}.\\
29/09 & Fernando Zegarra & \texttt{closure}, \texttt{goto}, \texttt{coleccionCanonica} y \texttt{llenarTablas}; reporte de conflictos y código de salida; \texttt{Makefile}; batería de 8 gramáticas con tabla esperada.\\
\bottomrule
\end{tabularx}
}

\medskip
\textbf{Decisiones de diseño}
\begin{itemize}[nosep,leftmargin=*]
\item \textbf{Ítem como tres enteros} (producción, punto, anticipación) y cada
      $cc_i$ como \texttt{std::set<Item>}: comparar dos conjuntos para saber si
      un $cc$ ya existe es directo y el orden de los ítems es determinista.
\item \textbf{Orden de las transiciones}: se precalcula el orden de primera
      aparición de los símbolos en la gramática (sin el inicial). Fija la
      numeración de los $cc_i$, que así coincide con la de otros equipos.
\item \textbf{Aumentación solo si hace falta}: si el símbolo inicial tiene una
      única producción y no aparece en ningún lado derecho (\texttt{Goal -> List}),
      esa regla hace de aumentada; si no, se agrega \texttt{S' -> S}.
\item \textbf{Conflictos}: cada celda de Action guarda sus acciones
      (\texttt{s3}, \texttt{r2}, \texttt{acc}) con los ítems que las generan.
      Si hay dos acciones distintas la herramienta no elige: reporta estado,
      terminal e ítems y sale con código 1.
\item \textbf{Código simple}: acciones como texto en vez de estructuras, una
      sola función para imprimir y comparar las tablas, y lista fija de casos
      de prueba (sin \texttt{<filesystem>}).
\item \texttt{first()} y \texttt{firstCadena()} se reutilizan sin cambios:
      \texttt{closure} calcula FIRST$(\delta a)$ llamando a \texttt{firstCadena}.
\end{itemize}

\medskip
\textbf{Problemas encontrados y cómo se resolvieron}
\begin{itemize}[nosep,leftmargin=*]
\item La carga original \emph{siempre} agregaba \texttt{S' -> S}. Con la
      gramática del enunciado eso creaba \texttt{Goal' -> Goal}: salían 13
      estados, la numeración se corría una posición y la tabla de referencia no
      coincidía. Se cambió a aumentación condicional (ver arriba).
\item \texttt{goto} es palabra reservada de C++: el método se llama
      \texttt{computeGoto}.
\item La producción vacía se guarda como \texttt{\{"epsilon"\}}: se añadió
      \texttt{body()}, que la devuelve vacía, para que su ítem sea directamente
      $[A \to \bullet,\, a]$ (completo).
\item El runner de pruebas recibía un único archivo; ahora recibe la carpeta
      de casos y compara cada gramática con su \texttt{.tabla}.
\end{itemize}

\medskip
\textbf{Verificación}
\begin{itemize}[nosep,leftmargin=*]
\item Caso de referencia (paréntesis): 12 conjuntos y la tabla exacta del enunciado.
\item Batería: \texttt{firstCadena} y 8 gramáticas con su tabla esperada, todas
      correctas. Incluye shift/reduce (\texttt{E -> E + E | id}), reduce/reduce
      (\texttt{A -> a}, \texttt{B -> a}) y producciones vacías.
\item Las tablas esperadas se generaron con una \emph{implementación
      independiente} en Python y la prueba falla si se altera una sola celda.
\end{itemize}

\medskip
\textbf{Pendiente}
\begin{itemize}[nosep,leftmargin=*]
\item Probar \texttt{make} y \texttt{make test} con g++: en la máquina de
      desarrollo solo había MSVC, con el que se compiló y probó todo en
      \texttt{/W4} sin avisos.
\end{itemize}

\end{document}
